\documentclass[10pt,a4paper]{article}
\usepackage[utf8]{inputenc}
\usepackage{geometry}
\usepackage{xcolor}
% Ajuste de margens mais estreitas para o layout de design
\geometry{a4paper, margin=0.75in}

% Definição de uma cor sobre para detalhes (ex: azul escuro/cinza)
\definecolor{primarycolor}{RGB}{30, 55, 90}

\begin{document}

% Cabeçalho de Destaque
\begin{center}
    {\Huge \textbf{\color{primarycolor}<< dados.basics.name >>}} \\
    \vspace{1mm}
    {\large \textit{<< dados.basics.label_pt if lang == 'pt' else dados.basics.label_en >>}} \\
    \vspace{2mm}
    \small << dados.basics.location >> | << dados.basics.phone >> | << dados.basics.email >>
\end{center}

\vspace{4mm}

% Início do Layout de Duas Colunas usando minipage
% Coluna Esquerda: Informações Técnicas e Contactos (33% da largura)
\begin{minipage}[t]{0.33\textwidth}
    \raggedright
    
    {\color{primarycolor}\large \textbf{<< 'Competências' if lang == 'pt' else 'Skills' >>}} \\
    \rule{\linewidth}{0.5pt} \\
    \vspace{1mm}
    % Loop para listar competências técnicas em linhas separadas
    <BLOCK> for skill in dados.skills.technical </BLOCK>
        • << skill >> \\
    <BLOCK> endfor </BLOCK>
    
    \vspace{5mm}
    
    {\color{primarycolor}\large \textbf{<< 'Idiomas' if lang == 'pt' else 'Languages' >>}} \\
    \rule{\linewidth}{0.5pt} \\
    \vspace{1mm}
    <BLOCK> for idioma in dados.skills.languages </BLOCK>
        \textbf{<< idioma.language_pt if lang == 'pt' else idioma.language_en >>} \\
        \small << idioma.fluency_pt if lang == 'pt' else idioma.fluency_en >> \\
        \vspace{1mm}
    <BLOCK> endfor </BLOCK>
\end{minipage}
\hfill
% Divisória invisível ou espaço entre colunas
\begin{minipage}[t]{0.02\textwidth}
    \hfill
\end{minipage}
\hfill
% Coluna Direita: Percurso Profissional e Académico (63% da largura)
\begin{minipage}[t]{0.63\textwidth}
    \raggedright
    
    {\color{primarycolor}\large \textbf{<< 'Experiência' if lang == 'pt' else 'Experience' >>}} \\
    \rule{\linewidth}{0.5pt} \\
    \vspace{2mm}
    
    <BLOCK> for xp in dados.experience </BLOCK>
        \textbf{<< xp.company >>} \hfill \small << xp.startDate >> -- << xp.endDate >> \\
        \textit{\small << xp.position_pt if lang == 'pt' else xp.position_en >>}
        \begin{itemize}
            <BLOCK> set destaques = xp.highlights_pt if lang == 'pt' else xp.highlights_en </BLOCK>
            <BLOCK> for item in destaques </BLOCK>
                \item \small << item >>
            <BLOCK> endfor </BLOCK>
        \end{itemize}
        \vspace{2mm}
    <BLOCK> endfor </BLOCK>
    
    \vspace{3mm}
    
    {\color{primarycolor}\large \textbf{<< 'Educação' if lang == 'pt' else 'Education' >>}} \\
    \rule{\linewidth}{0.5pt} \\
    \vspace{2mm}
    
    <BLOCK> for edu in dados.education </BLOCK>
        \textbf{<< edu.institution >>} \hfill \small << edu.startDate >> -- << edu.endDate >> \\
        \textit{\small << edu.area_pt if lang == 'pt' else edu.area_en >>} \\
        \small << edu.status_pt if lang == 'pt' else edu.status_en >> \\
        \vspace{2mm}
    <BLOCK> endfor </BLOCK>
\end{minipage}

\end{document}